\chapter{AV Fistula Creation}

\section*{What Is This Surgery?}
Arteriovenous (AV) fistula creation is a surgical procedure aimed at establishing a durable and high-flow vascular access for patients requiring long-term hemodialysis due to end-stage renal disease (ESRD). This operation involves the surgical connection of a native artery to a native vein, typically in the upper extremity, to facilitate direct arteriovenous communication. The primary goal of this procedure is to induce "arterialization" of the vein, leading to its dilation, thickening, and maturation into a robust conduit capable of withstanding repeated needle punctures necessary for hemodialysis sessions. AV fistulas are regarded as the gold standard for permanent vascular access due to their superior long-term patency, lower rates of infection, and reduced overall complications compared to synthetic grafts or central venous catheters. This surgical intervention significantly enhances patient quality of life and survival rates, making it a critical procedure in the management of patients with ESRD.

\section*{Alternative Names}
\begin{itemize}
  \item Arteriovenous fistula (AVF) creation
  \item Native fistula formation
  \item Cimino fistula
  \item Brescia-Cimino fistula
  \item Dialysis access surgery
  \item Vascular access creation
\end{itemize}

\section*{Relevant Surgical Anatomy}
The creation of an AV fistula primarily involves the vascular structures of the upper extremity, particularly the arteries and veins. Key arterial structures include the radial artery, ulnar artery, and brachial artery. These arteries provide the necessary high-pressure inflow essential for fistula maturation. The most commonly utilized veins for fistula creation are the cephalic vein, which runs superficially along the lateral aspect of the forearm and arm, and the basilic vein, which runs along the medial aspect of the forearm and arm, diving deep in the upper arm to join the brachial veins. The median cubital vein, connecting the cephalic and basilic veins in the antecubital fossa, is also a frequent site for anastomosis.

Surgical landmarks critical for AV fistula creation include the wrist for radiocephalic fistulas, the antecubital fossa for brachiocephalic or median cubital fistulas, and the upper arm for brachiobasilic fistulas, which may require transposition. Important nerve relations include the superficial radial nerve near the radial artery at the wrist, the median nerve in the antecubital fossa, and the ulnar nerve along the medial aspect of the upper arm. Careful dissection is essential to avoid nerve injury. The lymphatic drainage of the upper extremity generally follows the venous system, with superficial lymphatics draining into axillary nodes; surgical dissection should minimize lymphatic disruption to prevent lymphedema. Adequate arterial inflow, patent outflow veins without central stenosis, and suitable vessel diameters are critical for successful fistula formation and maturation.

\section*{Surgical Principle \& Rationale}
The fundamental principle underlying AV fistula creation is to establish a direct, low-resistance communication between a high-pressure arterial system and a low-pressure venous system. This direct anastomosis subjects the chosen vein to significantly increased arterial pressure and shear stress, initiating a biological process known as venous arterialization. Over weeks to months, the vein undergoes adaptive remodeling characterized by hypertrophy of its muscular wall, endothelial proliferation, and progressive dilatation. This transformation converts a thin-walled, collapsible vein into a robust, thick-walled conduit capable of withstanding repeated needle punctures and providing the high blood flow rates (typically >600 mL/min) required for efficient hemodialysis.

The rationale for creating a native AV fistula as the preferred method for permanent vascular access is multifaceted. It aims to provide a durable, self-healing access point that minimizes the risks associated with prosthetic materials (e.g., synthetic grafts) or temporary devices (e.g., central venous catheters). Key technical goals include achieving a widely patent, tension-free anastomosis, ensuring adequate arterial inflow and unobstructed venous outflow, and selecting vessels that possess the intrinsic capacity to mature to an appropriate diameter (typically >6 mm) and depth (<6 mm from the skin surface) over a sufficient length to facilitate easy and safe cannulation by dialysis personnel. The success of the fistula hinges on both meticulous surgical technique and the patient's inherent vascular health and adaptive response.

\section*{History \& Surgical Evolution}
The quest for reliable vascular access for hemodialysis began in the early 1960s. Initial attempts involved external shunts, most notably the Quinton-Scribner shunt, developed in 1960 by Belding Scribner and Wayne Quinton. This device connected an artery and vein using Teflon and silicone tubing, allowing for repeated access. While revolutionary for its time, these external shunts were plagued by high rates of infection and thrombosis, necessitating frequent revisions and limiting their long-term utility.

A pivotal breakthrough occurred in 1966 when Dr. James Cimino and Dr. Michael Brescia, working at the Bronx Veterans Administration Hospital, described the creation of a surgically constructed internal arteriovenous fistula in the forearm. This direct, subcutaneous connection between the radial artery and cephalic vein, now universally known as the Brescia-Cimino fistula, revolutionized hemodialysis access. Their innovation provided a durable, low-complication access that could be used for years, dramatically improving the quality of life and survival for patients with ESRD. Since this seminal work, various anatomical configurations (e.g., brachiocephalic, brachiobasilic) and refined surgical techniques have evolved, including the use of microsurgical instruments and magnification, but the underlying principle of creating a native AVF remains the preferred method for permanent hemodialysis access worldwide.

\section*{Clinical Indications}
AV fistula creation is indicated for patients with end-stage renal disease (ESRD) who require long-term hemodialysis. Specific clinical indications include:

\begin{itemize}
  \item To establish permanent vascular access for chronic hemodialysis in patients with ESRD, ideally before the anticipated need for dialysis.
  \item When a patient's peripheral veins are deemed inadequate for repeated cannulation, making other forms of access impractical or high-risk.
  \item To allow sufficient time for the fistula to mature (typically 4-12 weeks) before the anticipated initiation of maintenance hemodialysis.
  \item To reduce the risk of infection and thrombosis associated with temporary central venous catheters or synthetic arteriovenous grafts.
  \item As the preferred first-line option for vascular access due to its superior long-term patency and lower complication rates.
  \item For patients who have suitable native vessels identified during comprehensive preoperative mapping studies.
  \item In patients with progressive chronic kidney disease (CKD) Stage 4 or 5, where the glomerular filtration rate (GFR) is consistently below 20-25 mL/min, indicating impending need for renal replacement therapy.
\end{itemize}

\section*{Clinical Positioning}
AV fistula creation is overwhelmingly considered the primary, first-line procedure for establishing permanent vascular access in patients with end-stage renal disease requiring hemodialysis. Clinical guidelines, such as those from the Kidney Disease Outcomes Quality Initiative (KDOQI), strongly advocate for a native arteriovenous fistula as the initial choice due to its superior long-term patency, significantly lower rates of infection, and reduced need for reintervention compared to synthetic grafts or central venous catheters. It is typically performed several months (ideally 3-6 months) before the anticipated initiation of dialysis to allow for adequate maturation.

In certain circumstances, an AV fistula might be considered a secondary or salvage option. This includes situations where previous AVF attempts have failed, when suitable native vessels are not available in the distal limb, or if specific patient comorbidities (e.g., severe peripheral arterial disease, central venous stenosis) preclude a more distal or primary fistula. Even in these scenarios, the goal remains to create a native fistula, often a more proximal one (e.g., brachiocephalic or brachiobasilic) or one requiring vein transposition, whenever anatomically and clinically feasible, before resorting to prosthetic grafts or permanent catheters.

\section*{Surgical Technique \& Procedure}
The creation of an AV fistula is a precise surgical procedure typically performed under regional anesthesia (e.g., axillary block) or local anesthesia with sedation, though general anesthesia may be utilized in select cases. The patient is positioned supine with the chosen arm abducted and externally rotated on an arm board, ensuring optimal exposure and comfort.

The key steps of the procedure generally include:

\begin{itemize}
  \item \textbf{1. Anesthesia and Patient Positioning:} After appropriate anesthesia, the chosen limb is prepped and draped in a sterile fashion. The arm is positioned to allow full access to the planned incision site and to ensure the vessels are not under tension.
  
  \item \textbf{2. Incision and Dissection:} A small transverse or longitudinal incision, typically 3-5 cm in length, is made over the planned anastomosis site (e.g., radial artery and cephalic vein at the wrist, or brachial artery and cephalic/basilic vein at the elbow). Meticulous dissection is performed to identify and carefully mobilize the target artery and vein, preserving surrounding nerves, lymphatic structures, and collateral vessels.

  \item \textbf{3. Vessel Preparation:} The artery and vein are gently mobilized, and their branches are ligated or clipped. Vessel loops or soft clamps are used for proximal and distal control. The vessels are then flushed with heparinized saline to prevent thrombus formation, and systemic heparinization may be administered to reduce intraoperative clotting.

  \item \textbf{4. Anastomosis Creation:} The chosen artery and vein are prepared for anastomosis. This typically involves a venotomy (incision into the vein) and an arteriotomy (incision into the artery). The most common configurations are end-to-side (vein end to artery side), side-to-side, or end-to-end, depending on vessel anatomy and surgeon preference. The anastomosis is meticulously performed using fine, non-absorbable monofilament sutures (e.g., 6-0 or 7-0 Prolene) under magnification, ensuring a wide, tension-free, and leak-proof connection.

  \item \textbf{5. Confirmation of Patency:} After completing the anastomosis and removing the vessel clamps, blood flow is restored. The surgeon immediately confirms the patency of the fistula by palpating a strong thrill (a characteristic vibration) and auscultating a continuous bruit (a whooshing sound) over the anastomosis and along the course of the vein. Any signs of inadequate flow or thrill warrant immediate re-evaluation.

  \item \textbf{6. Hemostasis and Closure:} Meticulous hemostasis is achieved by controlling any bleeding points. The wound is then closed in layers, typically with absorbable sutures for subcutaneous tissues and non-absorbable sutures or staples for the skin. Drains are rarely used. A sterile dressing is applied, often with light compression.
\end{itemize}

The specific type of fistula (e.g., radiocephalic, brachiocephalic, brachiobasilic with or without transposition) is determined preoperatively based on vein mapping and patient anatomy, always prioritizing the most distal and native vessel option.

\section*{Preoperative Workup \& Preparation}
Thorough preoperative workup and preparation are paramount for optimizing the success of AV fistula creation and maturation. This comprehensive process typically involves:

\begin{itemize}
  \item \textbf{Clinical Assessment:} A detailed medical history and physical examination, focusing on cardiovascular status, peripheral pulses, and a meticulous evaluation of both upper extremities for suitable vessels. Assessment for any signs of infection or peripheral arterial disease is crucial.

  \item \textbf{Preoperative Imaging (Vein Mapping):} Duplex ultrasound is the gold standard for mapping the arterial and venous anatomy of the chosen limb. This study identifies vessel patency, diameter (arterial diameter ideally >2 mm, venous diameter >2.5 mm for wrist, >3 mm for elbow), depth, and the presence of any stenosis, thrombus, or calcification, guiding the surgeon in selecting the optimal site and type of fistula.

  \item \textbf{Laboratory Tests:} Routine blood tests including a complete blood count, coagulation profile (PT/INR, aPTT), renal function tests (creatinine, BUN), and electrolytes are performed to assess overall health, identify any bleeding risks, and confirm ESRD status.

  \item \textbf{Medication Management:} Patients on anticoagulants or antiplatelet agents (e.g., warfarin, clopidogrel, aspirin) may require temporary cessation or bridging strategies as per institutional protocols, balancing the risk of bleeding with the risk of thrombosis. Essential medications should be continued.

  \item \textbf{NPO Status:} Patients are typically instructed to be NPO (nil per os) for a specified period before surgery, usually 6-8 hours, to minimize the risk of aspiration during anesthesia.

  \item \textbf{Prophylactic Antibiotics:} A single dose of intravenous prophylactic antibiotics (e.g., cefazolin) is often administered within 60 minutes prior to skin incision to reduce the risk of surgical site infection.

  \item \textbf{Informed Consent:} A detailed discussion with the patient regarding the procedure, its benefits, potential risks, alternative access options, and the critical importance of postoperative care and the maturation period.

  \item \textbf{Arm Preservation:} Patients are educated to strictly avoid venipuncture, intravenous lines, and blood pressure measurements in the limb designated for fistula creation, both pre- and postoperatively, to preserve venous integrity and patency.
\end{itemize}

Ideally, fistula creation should occur several months (at least 2-3 months) before the anticipated need for hemodialysis to allow adequate time for maturation.

\section*{Contraindications \& Precautions}
\textbf{Absolute Contraindications:}
\begin{itemize}
  \item Severe peripheral arterial disease (PAD) in the limb designated for access, leading to inadequate arterial inflow and high risk of steal syndrome or fistula failure.
  \item Central venous stenosis or occlusion in the ipsilateral limb, which would impede venous outflow, prevent maturation, and lead to venous hypertension.
  \item Active systemic infection or localized infection in the limb designated for fistula creation.
  \item Uncorrectable coagulopathy or severe bleeding disorder that cannot be managed perioperatively.
  \item Severe, uncontrolled congestive heart failure (NYHA Class IV), where the increased cardiac output demand from the fistula could acutely exacerbate heart failure.
  \item Unsuitable native vessels on preoperative duplex ultrasound mapping (e.g., excessively small diameter, thrombosed, sclerosed, or deeply located veins).
\end{itemize}

\textbf{Relative Contraindications \& Precautions:}
\begin{itemize}
  \item Severe cardiac dysfunction (e.g., ejection fraction <30\%), where the risk of high-output cardiac failure is elevated, requiring careful consideration and cardiology clearance.
  \item Morbid obesity, which can make vessels deep and challenging to identify, dissect, and cannulate, potentially increasing the risk of maturation failure.
  \item Uncontrolled diabetes mellitus with severe microvascular disease, potentially affecting vessel quality, healing, and increasing the risk of steal syndrome.
  \item Previous failed access attempts in the same limb, indicating potentially compromised vasculature and requiring thorough re-evaluation.
  \item Patient non-compliance or inability to understand and adhere to postoperative care instructions, which is crucial for maturation and long-term patency.
  \item Preserving veins in the non-access arm for future potential access or other critical medical needs (e.g., chemotherapy, IV antibiotics).
  \item Avoiding blood draws, intravenous lines, and blood pressure cuffs on the access arm both pre- and postoperatively to protect the developing fistula.
\end{itemize}

\section*{Complications \& Risks}
While AV fistula creation is generally safe and effective, several potential risks and complications are associated with the procedure:

\begin{itemize}
  \item \textbf{General Surgical Complications:}
  \begin{itemize}
    \item \textit{Bleeding and Hematoma:} Accumulation of blood at the surgical site, which can cause pain, swelling, and rarely require re-exploration.
    \item \textit{Infection:} Surgical site infection (cellulitis, abscess), or, rarely, systemic sepsis, particularly if the fistula becomes infected after cannulation.
    \item \textit{Anesthesia Complications:} Risks associated with local, regional, or general anesthesia, including allergic reactions, respiratory depression, cardiac events, or nerve injury from block placement.
  \end{itemize}
  
  \item \textbf{Procedure-Specific Risks:}
  \begin{itemize}
    \item \textit{Maturation Failure:} The most common complication, where the vein fails to dilate and thicken sufficiently (e.g., to the "Rule of 6s" criteria) for successful cannulation and effective dialysis. This can be due to inadequate flow, venous stenosis, or poor vessel quality.
    \item \textit{Thrombosis:} Formation of a blood clot within the fistula, leading to occlusion and loss of function. This can occur early postoperatively or years later, often requiring urgent intervention.
    \item \textit{Stenosis:} Narrowing of the vein, often at the anastomosis, in the outflow vein, or in central veins, which can reduce flow, impede maturation, and lead to thrombosis or venous hypertension.
    \item \textit{Peripheral Ischemia (Steal Syndrome):} Diversion of arterial blood flow away from the hand or fingers, leading to symptoms such as pain, numbness, coldness, pallor, or even tissue loss (ulceration, gangrene) in severe cases.
    \item \textit{Aneurysm or Pseudoaneurysm Formation:} Localized dilatation of the vein (aneurysm) or a contained rupture (pseudoaneurysm) of the vessel wall, often due to repeated cannulation trauma, high flow, or underlying vessel wall weakness.
    \item \textit{High-Output Cardiac Failure:} In rare cases, a very high-flow fistula can significantly increase cardiac preload and output, potentially exacerbating or precipitating heart failure in susceptible individuals with pre-existing cardiac disease.
    \item \textit{Venous Hypertension:} Swelling, pain, and discoloration of the hand or forearm distal to the fistula due to increased venous pressure and impaired drainage, often caused by outflow stenosis.
    \item \textit{Nerve Injury:} Damage to adjacent nerves (e.g., median, ulnar, radial nerves) during surgical dissection, leading to numbness, weakness, paresthesia, or pain in the distribution of the affected nerve.
    \item \textit{Failure to Cannulate:} Even if technically mature, some fistulas may be difficult to cannulate due to excessive depth, tortuosity, or patient factors, leading to prolonged reliance on catheters.
  \end{itemize}
\end{itemize}

\section*{Postoperative Recovery Timeline}
Immediately after AV fistula creation, the patient is transferred to the post-anesthesia care unit (PACU) for close monitoring of vital signs, pain control, and initial assessment of the access site. The surgeon or nursing staff will confirm the presence of a palpable thrill and an audible bruit over the fistula, indicating successful blood flow. The arm is typically elevated to reduce swelling, and appropriate pain medication is administered.

Over the first 24-48 hours, patients are usually discharged home with detailed instructions for wound care, pain management, and activity restrictions. They are advised to avoid heavy lifting, direct pressure on the fistula, and tight clothing or jewelry on the access arm. Crucially, patients are taught to regularly check for the thrill and bruit, as their absence can indicate thrombosis and requires immediate medical attention. The most critical phase is the maturation period, which typically lasts 4 to 12 weeks, but can extend longer. During this time, the vein gradually dilates and thickens, becoming suitable for cannulation. Gentle exercises, such as squeezing a rubber ball, may be encouraged to promote venous development. Long-term recovery involves vigilant self-monitoring, adherence to arm preservation rules, and regular follow-up appointments to monitor maturation and address any emerging complications, ensuring the fistula is ready for sustained hemodialysis.

\section*{Surgical Findings \& Pathology}
During AV fistula creation, the surgeon's findings are critical for assessing the suitability of vessels and the immediate success of the procedure. Key intraoperative findings include the quality and diameter of the chosen artery (e.g., presence of calcification, adequate pulsatility) and vein (e.g., wall thickness, presence of valves, absence of thrombus or significant scarring). After anastomosis, the immediate presence of a strong, continuous thrill and a distinct bruit confirms patency and adequate flow. The surgeon also assesses for any signs of tension at the anastomosis or kinking of the vessels.

While routine pathology on resected tissue is uncommon for AV fistula creation, if a segment of vein or artery is excised (e.g., for revision or if a segment is deemed unsuitable), the pathology report would typically describe the vessel wall architecture. In the context of a maturing fistula, the vein wall would show signs of arterialization, including medial hypertrophy (thickening of the smooth muscle layer), intimal hyperplasia (thickening of the innermost layer), and elastic tissue remodeling. These changes are adaptive responses to the increased pressure and shear stress, transforming the vein into a more robust, artery-like conduit. In cases of failed maturation or complications, pathology might reveal excessive intimal hyperplasia leading to stenosis, thrombosis, or signs of infection.

\section*{Expected Postoperative Course}
A normal and successful postoperative course following AV fistula creation involves a period of healing and maturation, culminating in a functional access for hemodialysis. Initially, patients can expect mild to moderate pain at the incision site, managed with oral analgesics, and some swelling in the hand and forearm, which typically subsides within a few days to weeks. The surgical wound should heal cleanly, with sutures or staples removed around 10-14 days post-op.

The most crucial aspect of the expected course is the maturation of the fistula. Over 4 to 12 weeks, the vein should progressively dilate, its wall should thicken, and blood flow should increase. Patients should regularly feel a strong, continuous thrill and hear a distinct bruit, which are signs of a healthy, patent fistula. The vein should become visibly prominent and palpable, yet soft and compressible, making it suitable for cannulation. A successful course means the fistula meets the "Rule of 6s" criteria (diameter >6mm, flow >600mL/min, depth <6mm, usable length >6cm) and can be reliably cannulated for efficient hemodialysis without complications like steal syndrome, infection, or thrombosis.

\section*{Postoperative Care \& Follow-up}
Postoperative care and diligent follow-up are critical for ensuring AV fistula maturation and long-term patency.

\begin{itemize}
  \item \textbf{Initial Postoperative Visits:} Typically scheduled 1-2 weeks after surgery for wound check, suture/staple removal, and initial assessment of thrill and bruit.
  
  \item \textbf{Maturation Monitoring:} Regular clinical assessments (physical examination for thrill, bruit, swelling, and signs of maturation) are performed every 2-4 weeks during the maturation period (4-12 weeks).
  
  \item \textbf{Duplex Ultrasound Surveillance:} Serial duplex ultrasound scans may be performed to objectively assess flow rates, vessel diameter, depth, and to identify any developing stenosis or other issues that might impede maturation.
  
  \item \textbf{Patient Education:} Reinforcement of self-monitoring techniques (checking for thrill/bruit daily), instructions on arm care (avoiding tight clothing, heavy lifting, sleeping on the arm), and strict adherence to arm preservation rules (no blood draws, IVs, or BP cuffs on the access arm).
  
  \item \textbf{Activity Resumption:} Gradual return to normal activities, avoiding strenuous activities or heavy lifting with the access arm until full maturation is confirmed.
  
  \item \textbf{Physical Therapy/Exercises:} Gentle hand and forearm exercises (e.g., squeezing a soft rubber ball) may be recommended to promote venous dilatation and strengthening.
  
  \item \textbf{Warning Signs:} Patients are instructed to contact their healthcare provider immediately if they experience loss of thrill/bruit, sudden swelling, severe pain, redness, warmth, pus discharge, numbness, or coldness in the access arm or hand.
\end{itemize}

If maturation failure or complications like stenosis or thrombosis occur, timely interventions such as angioplasty, stenting, or surgical revision may be required to salvage the fistula.

\section*{Epidemiology}
Arteriovenous fistula creation is one of the most frequently performed vascular surgical procedures globally, driven by the increasing prevalence and incidence of end-stage renal disease (ESRD). In the United States, over 800,000 individuals are living with ESRD, with the majority requiring hemodialysis. Annually, tens of thousands of AV fistulas are created, reflecting the continuous need for reliable vascular access. The incidence of ESRD is higher in older adults, with a median age of diagnosis typically in the mid-60s, and while there isn't a strong sex predilection for ESRD itself, women tend to have smaller native vessels, which can sometimes pose challenges for fistula creation and maturation, leading to higher rates of early failure. The primary indication for AVF creation is ESRD requiring chronic hemodialysis, accounting for nearly all procedures. While the mortality directly associated with AVF creation is very low, the morbidity related to maturation failure (up to 20-50\%), thrombosis, and subsequent reinterventions is significant, impacting patient quality of life and imposing substantial healthcare costs.

\section*{Outcomes \& Prognosis}
The outcomes and prognosis for AV fistula creation are generally favorable, making it the preferred method for hemodialysis access. Primary patency rates (fistula remaining functional without intervention) typically range from 60-80\% at one year and 40-70\% at three years, depending on patient factors, fistula type, and surgical technique. Secondary patency rates (fistula remaining functional with interventions to maintain patency) are even higher, often exceeding 80-90\% at one year and 70-80\% at three years, highlighting the importance of surveillance and timely salvage procedures. Patients with a functioning AVF experience significantly lower rates of infection, thrombosis, and hospitalization compared to those with synthetic grafts or central venous catheters, leading to improved survival and quality of life. Prognostic factors influencing outcomes include patient comorbidities (e.g., diabetes, peripheral vascular disease, obesity), age, vessel quality (diameter, calcification, presence of central stenosis), and the experience of the surgical team. While maturation failure remains a common early complication, successful salvage procedures can often convert a failing fistula into a functional one, underscoring the importance of vigilant surveillance and timely intervention, as emphasized by guidelines such as the KDOQI.

\section*{References}
\begin{enumerate}
  \item National Kidney Foundation. KDOQI Clinical Practice Guideline for Vascular Access: 2019 Update. American Journal of Kidney Diseases. 2019;74(4 Suppl 2):S1-S150.
  
  \item Sidawy AN, Perler BA. Rutherford's Vascular Surgery and Endovascular Therapy. 10th ed. Elsevier; 2023.
  
  \item MedlinePlus. Hemodialysis access. U.S. National Library of Medicine; 2024.
  
  \item Sabiston DC, Townsend CM. Sabiston Textbook of Surgery: The Biological Basis of Modern Surgical Practice. 21st ed. Elsevier; 2021.
\end{enumerate}

\clearpage